\section{Introduction}
\label{sec:introduction}

\subsection{Motivation}

Accurately forecasting the magnitude of future asset returns is a central problem in
quantitative finance. While the \emph{direction} of price changes is notoriously
difficult to predict, the \emph{amplitude} of those changes (often proxied by realized
volatility) exhibits strong and exploitable serial dependence. This phenomenon is
characterised by \emph{volatility clustering}: large
price moves tend to be followed by large moves, and calm periods persist until a new
shock arrives.

Predicting the absolute return of the S\&P~500 index therefore serves a dual purpose: it
is both a well-posed statistical problem with a meaningful signal in the data and a
practically relevant task for risk management, option pricing, and portfolio construction.

\subsection{Objective}

This project empirically compares a progression of models of increasing expressive power
applied to the same forecasting task:

\begin{enumerate}
  \item \textbf{Linear Regression} --- a transparent but structurally constrained baseline.
  \item \textbf{GARCH(1,1)} --- the industry-standard parametric volatility model.
  \item \textbf{Feedforward Neural Network (FNN)} --- a non-linear model that treats the
        input as a flat, orderless feature vector.
  \item \textbf{Recurrent Neural Network (RNN)} --- a sequential model that can exploit
        the temporal ordering of lags.
  \item \textbf{Long Short-Term Memory (LSTM)} --- a gated architecture designed to
        capture long-range dependencies.
\end{enumerate}

Rather than searching for the best possible model, the goal is to \emph{understand the
limitations of each architecture} and to show how each successive model overcomes
(some of) the shortcomings of its predecessor.

\subsection{Dataset}

The dataset is built from the daily log-returns of the SPY ETF (S\&P~500 tracker),
downloaded via the \texttt{yfinance} library. The sample spans \textbf{February 1993 to
March 2026}, yielding \textbf{8\,343 daily observations}.

The forecasting problem is formulated as a sliding-window regression:
\begin{itemize}
  \item \textbf{Input} $X_t$: the $L$ most recent daily log-returns
        $(r_{t-1}, r_{t-2}, \ldots, r_{t-L})$, with $L \geq 10$.
  \item \textbf{Target} $y_t$: the realized absolute return on day $t$, computed as the
        rolling standard deviation of returns over a short window (a proxy for
        next-day volatility).
\end{itemize}

\subsection{Evaluation}

Models are evaluated using three metrics, computed on a held-out test set:

\begin{itemize}
  \item \textbf{MAE} (Mean Absolute Error) --- primary metric for cross-model comparison.
  \item \textbf{RMSE} (Root Mean Squared Error) --- penalises large errors more heavily.
  \item \textbf{QLIKE} (Quasi-Likelihood loss) --- a standard volatility-forecasting loss
        that penalises systematic under-prediction.
\end{itemize}

Because market conditions vary dramatically, all models are also evaluated separately on
two \textbf{market regimes} identified from the data:
\begin{itemize}
  \item \textbf{Calm regime} ($\approx$90\% of observations): low, stable volatility.
  \item \textbf{Crisis regime} ($\approx$10\% of observations): sharp spikes corresponding
        to the dot-com crash, the 2008 Global Financial Crisis, and the COVID shock.
\end{itemize}

The cross-model summary (Section~\ref{sec:lstm}) reports MAE only.

\subsection{Paper Structure}

Section~\ref{sec:data} describes the dataset and its key statistical properties.
Section~\ref{sec:linreg} presents the linear regression and GARCH baselines.
Sections~\ref{sec:fnn}--\ref{sec:lstm} progressively introduce the neural network
architectures. A final comparison table closes the report.
